% No modificar estas líneas de código, por favor dirigirse a CONCLUSIONES

\newpage
\phantomsection\addcontentsline{toc}{section}{CONCLUSIONES} 

%------------------------
%
%       CONCLUSIONES
%
%------------------------

\section*{CONCLUSIONES}
Introducir lo que se desea....ahora aprovecharemos para demostrar como adicionar las bibliografías. 

Al igual que ocurre con las ecuaciones, tablas y figuras, las referencias de la bibliografía se pueden citar en el texto mediante la etiqueta que le pongamos a cada una dentro de \verb|bibliografia.bib|.

Por ejemplo, si quiero referirme al artículo, libro, página web, etc. lo puedo hacer escribiendo \verb|\citep{etiqueta}|.

Si la referencia que ponemos en la bibliografía es de internet, se debe adicionar la dirección web para que pueda ser consultada y la fecha en que se ingresó. Por supuesto, los enlaces que se pongan deben dirigir a páginas que cumplan con todos los requisitos legales en cuanto al respeto a los derechos de la propiedad intelectual \citep{pressman1988ingenieria}.

Igual se cuenta con otras formas de citado, que se utiliza cuando se hablará de forma directa de un autor, es decir: "Según \cite{AI} blah blah....." se utilizará el comando \verb|\cite{etiqueta}|. 

Caso usted desee hacer una citación textual de algún autor o posiblemente entrevista, se debe utilizar el espacio de trabajo \verb|\begin{quote}, \end{quote}|. El mismo se presenta a continuación, destacando que se debe comenzar con la introducción del autor con el comando \verb|\cite{etiqueta}| o en su defecto al final de la citación con el comando \verb|\citep{etiqueta}|.

\begin{quote}
\singlespacing
\textit{"Blah blah blah blah blah blah blah blah blah blah blah blah blah blah blah blah blah blah blah blah blah blah blah blah blah blah blah blah blah blah blah blah blah blah blah blah."}    
\end{quote}

Entre otras variaciones contamos con el formato de citado Cf. con el comando \verb|\citecf{etiqueta}|, todas las formas de citación deben ser utilizadas con el cuidado correspondiente. 

